% ---------------------------------------------------------------------
%  CHƯƠNG 2 — CƠ SỞ LÝ THUYẾT VÀ CÔNG NGHỆ
% ---------------------------------------------------------------------
\chapter{Cơ sở lý thuyết và công nghệ}
\label{ch:coso}

\begin{muctieu}
Chương này tóm lược các kiến thức nền tảng của học phần được dùng trực tiếp
trong đồ án và danh mục công cụ phát triển.
\end{muctieu}

\section{Kiến thức nền tảng}

\begin{dinhnghia}[Trạng thái dùng chung]
Trạng thái dùng chung \en{(shared state)} là dữ liệu được nhiều luồng cùng đọc
và ghi. Trong đồ án này, danh sách đội và lịch sử trận đấu là trạng thái dùng
chung: mỗi trọng tài là một luồng ghi thêm một trận, mỗi khách xem là một luồng
đọc bảng xếp hạng.
\end{dinhnghia}

Đồ án dùng bốn mảng kiến thức chính của học phần. Thứ nhất, cấu trúc dữ liệu
chương~3: \code{Map} để tra cứu đội theo tên với độ phức tạp логарифmic, \code{List}
để lưu lịch sử trận theo thứ tự thời gian. Thứ hai, streams và generics
chương~2: thống kê bảng xếp hạng được tính bằng \code{stream} trên tập đội, sắp
xếp bằng \code{Comparator} có chuỗi \code{thenComparing}. Thứ ba, lập trình
đa luồng chương~4: \code{Thread}, \code{Runnable}, \code{synchronized},
\code{ExecutorService} và \code{CopyOnWriteArrayList}. Thứ tư, lập trình mạng
chương~5: \code{ServerSocket}, \code{Socket}, giao thức văn bản tự định nghĩa.
Các tính năng hiện đại của Java~6 được dùng xuyên suốt: \code{record} cho dữ
liệu bất biến, \code{enum} cho lệnh giao thức, \code{try-with-resources} khi mở
tệp và socket, và \code{switch} dạng biểu thức~\cite{Bloch2018}.

\begin{ghinho}
Nguyên tắc đồng bộ được chọn là: đồng bộ ở mức một thao tác nghiệp vụ (nhập một
trận) chứ không phải ở mức từng trường, đảm bảo nguyên tử mà vẫn để nhiều luồng
đọc bảng xếp hạng song song.
\end{ghinho}

\section{Công cụ và thư viện}

Đồ án chỉ dùng thư viện chuẩn của Java, không dùng Spring hay thư viện ngoài.
Bảng~\ref{tab:congcu} liệt kê các thành phần chính.

\begin{table}[H]
\centering
\caption{Công cụ và thư viện sử dụng trong đồ án.}
\label{tab:congcu}
\begin{tabular*}{\textwidth}{@{\extracolsep{\fill}}>{\bfseries\color{daunavy}}l >{\color{daumaroon}}l l}
\toprule
\rowcolor{gray100}\bfseries\color{daunavy}Thành phần & \bfseries\color{daumaroon}Phiên bản & \bfseries\color{daunavy}Vai trò \\
\midrule
JDK            & 17+ (thử trên 21) & Ngôn ngữ và thư viện chuẩn \\
java.net       & --   & ServerSocket, Socket TCP \\
java.util.concurrent & -- & ExecutorService, CopyOnWriteArrayList \\
java.nio.file  & --   & Đọc/ghi CSV bằng try-with-resources \\
git            & 2.x  & Quản lý phiên bản \\
PowerShell     & --   & Biên dịch và chạy trên Windows \\
\bottomrule
\end{tabular*}
\end{table}

\section{Các công trình liên quan}

Ví dụ máy chủ echo TCP đa luồng trong giáo trình chương~5 là cơ sở trực tiếp
của đồ án: cùng dùng \code{ServerSocket} và một luồng phục vụ mỗi máy khách.
Đồ án mở rộng ví dụ đó bằng cách thêm trạng thái dùng chung cần đồng bộ, giao
thức đa lệnh có đóng khung thông điệp, và tầng dữ liệu có lưu/khôi phục. Cách
dùng \code{CopyOnWriteArrayList} cho danh sách người nghe và khoá riêng khi
phát broadcast tham khảo cách tổ chức của các hệ thống pub/sub nhẹ trong
\cite{Goetz2006}.
